\chapter{Tracing and observability}
\label{ch:tracing}

\begin{aside}
You cannot fix what you cannot measure. The renderer's quiet
nature is also its trickiness: a slow frame, a sudden burst of
output, a memory leak --- nothing screams. Tracing is how you
find out.
\end{aside}

\section{Trace points}

A trace point is a recorded event with a name and a timestamp.
\maya{} sprinkles them through the render pipeline:

\begin{lstlisting}
trace("frame_begin");
layout();
trace("layout_done");
paint();
trace("paint_done");
diff();
trace("diff_done");
emit();
trace("frame_end");
\end{lstlisting}

Compiled out unless \code{MAYA\_TRACE} is set.

\section{Output format}

Traces are written as JSON, one event per line, compatible
with chrome://tracing:

\begin{lstlisting}
{"name":"layout","ph":"X","ts":12345,"dur":540,
 "pid":1,"tid":1}
\end{lstlisting}

Tools like Perfetto or Speedscope visualise the timeline.

\section{Per-frame counters}

Beyond timing, \maya{} records counters per frame:

\begin{itemize}[leftmargin=*]
  \item Cells changed.
  \item ANSI bytes emitted.
  \item Effects re-run.
  \item Elements in tree.
  \item Allocations.
\end{itemize}

A status bar can display these for debugging. \maya{} ships
\code{stats\_widget()} that does exactly this.

\section{Logging levels}

A simple log filter:

\begin{itemize}[leftmargin=*]
  \item \code{ERROR}: something is broken; user-visible.
  \item \code{WARN}: surprising but recoverable.
  \item \code{INFO}: notable events (startup, shutdown,
    capability detection).
  \item \code{DEBUG}: diagnostic detail.
  \item \code{TRACE}: per-frame, per-event spam.
\end{itemize}

The active level is environment-set: \code{MAYA\_LOG=DEBUG}.

\section{Where logs go}

Not to stderr (the terminal is the display, not a log sink).
\maya{} writes logs to a file: \code{maya.log} in the current
directory by default, configurable.

\begin{tryit}
Set \code{MAYA\_LOG=DEBUG} and tail the log file in another
terminal. Watch the runtime narrate itself.
\end{tryit}

\section{Counters and aggregates}

For long-running apps, summary stats every minute:

\begin{itemize}[leftmargin=*]
  \item Frames per second (during activity).
  \item Average ANSI bytes per frame.
  \item Peak memory.
  \item Total signal updates.
\end{itemize}

These accumulate cheaply (one integer increment per event)
and dump periodically.

\section{User-visible diagnostics}

Some libraries hide their internals; others expose them. \maya{}
exposes via:

\begin{itemize}[leftmargin=*]
  \item \code{maya::about()}: prints library version, build
    flags, detected terminal, capability summary.
  \item \code{maya::stats()}: dumps current counters.
  \item Debug overlay: live counters and rectangles.
\end{itemize}

Users hitting odd behaviour can paste the output of
\code{about()} in a bug report and the maintainer knows
what to look at first.

\section{Profiling integration}

For production profiling, \maya{} cooperates with:

\begin{itemize}[leftmargin=*]
  \item \code{perf}: standard Linux sampling.
  \item \code{tracy}: instrumentation-based.
  \item \code{Tracing Framework} or \code{TEF}: chrome
    JSON.
\end{itemize}

Trace points compile down to the chosen backend via macro
glue. Tracy gets native zones; perf gets nothing extra; TEF
gets JSON.

\section{Signal-graph observability}

The signal graph (see Chapter~\ref{ch:signals}) can dump its
state: which signals exist, which effects subscribe to what,
last update time. This is invaluable for debugging
``something updates every frame'' regressions.

\begin{lstlisting}
ctx.dump_graph("/tmp/graph.dot");
\end{lstlisting}

Pipe through \code{dot -Tpng} for a visual.

\section{Recap}

\begin{recap}
\begin{itemize}[leftmargin=*]
  \item Trace points compile out by default.
  \item JSON format is chrome-compatible.
  \item Per-frame counters in a stats widget.
  \item Logs to file, not stderr.
  \item \code{about()} for bug reports.
\end{itemize}
\end{recap}

\section*{Workshop}

\begin{enumerate}[leftmargin=*]
  \item Add the stats widget to your status bar. Identify a
    frame that took unusually long.
  \item Set \code{MAYA\_TRACE=1} and visualise a 5-second
    capture in chrome://tracing.
  \item Print \code{maya::about()} on startup. Confirm it
    reports your terminal and capabilities.
\end{enumerate}
